%**************************************************************************
%*                      Example of paper for the                          *
%*         ELECTRONIC LETTERS ON COMPUTER VISION AND IMAGE ANALYSIS       *
%*                      LaTex 2e Document                                 *
%**************************************************************************

\documentclass[twoside,a4paper,11pt]{article}
\usepackage{times}
\usepackage{elcvia}
\usepackage[english]{babel}

\newcommand{\Runo}{\rm I\hspace{-0.55ex}R}

\input{psfig.sty}
%\psdraft

% INFORMATION TO FILL IN CAMERA READY VERSION
%------------------------------------------------------------------
\newcommand{\vol}{0}            % Number of Volume (default 0)
\newcommand{\num}{0}            % Number (default 0)
\setcounter{page}{1}            % First page (default 1)
\newcommand{\begpag}{\thepage}  
\newcommand{\finpag}{7}         % Last Page
\newcommand{\yearp}{2000}       % Year of the issue (default 2000)
%------------------------------------------------------------------
\renewcommand{\thefootnote}{\fnsymbol{footnote}}

\title{Example of paper for the Electronic Letters on Computer Vision
and Image Analysis 
\footnote[0]{Correspondence to: $<$E-mail$>$ \vspace{0.2cm}}
% INFORMATION TO FILL IN CAMERA READY VERSION
%------------------------------------------------------------------
% name of Associate Editor given when article is accepted
\footnote[0]{\noindent Recommended for acceptance by $<$name$>$} 
%------------------------------------------------------------------
\footnote[0]{ELCVIA ISSN:1577-5097}
\footnote[0]{Published by Computer Vision Center / Universitat Aut\`onoma de 
Barcelona, Barcelona, Spain}
}

\author{Name1$^*$ and Name2$^+$ \\ \\
\centerline{\small \em $^*$ Department, Research Center, Address, City, Country} \\
\centerline{\small \em $^+$ Department, Research Center, Address, City, Country} \\ \\
% INFORMATION TO FILL IN CAMERA READY VERSION
%------------------------------------------------------------------
% Received, revised and accepted dates
\centerline{\small Received 1 January 2000; revised 1 January 2000; accepted 1 January 2000}
       }


\begin{document}

\maketitle

\pagestyle{myheadings}

% INFORMATION TO FILL IN CAMERA READY VERSION
%-------------------------------------------------------------------
% fill in (authors) the authors of article in following format:
% * 1,2 authors: J. Jones and J. Jones
% * More than 2 authors: J. Jones et al.
\markboth{\centerline{\small \it Name1 et al. / 
           Electronic Letters on Computer Vision and Image Analysis 
           \vol(\num):\begpag-\finpag, \yearp}}
         {\centerline{\small \it Name1 et al. / 
           Electronic Letters on Computer Vision and Image Analysis
           \vol(\num):\begpag-\finpag, \yearp}}
%-------------------------------------------------------------------
\hrulefill

\begin{abstract}

This paper wants to be a presentation of the e-journal
{\em Electronic Letters on Computer Vision and Image Analysis} (ELCVIA).

ELCVIA is a new electronic journal whose aim is to provide
a forum where Computer Vision and Image Analysis researchers
can discuss about methods, applications and experiences in this field.

ELCVIA tries to offer to authors electronical means (images, video or 
exe files) to show in a better form their results, which are
more difficult to demonstrate in paper format.

ELCVIA tries to make the usual journal process (article submission, referee's
revision and revisions reception) through WWW, by means of Web forms in the ELCVIA
Web site {\sf www.cvc.uab.es/elcvia}.

ELCVIA offers two types of contributions: DRAFT, with an open discussion of
all the people, and ARTICLE, with a close revision of Editorial Board and
reviewers. 

This paper also contains information to the authors about the article or draft 
style format to be send to ELCVIA.

This example is a {\bf \LaTeX2e version.}
 
\vspace{0.5cm}
\noindent
{\em Key Words}: Computer Vision, Image Analysis,
Pattern Recognition Image segmentation, 3D Reconstruction,
Active Vision, Tracking, Video and Image Sequence Analysis, \LaTeX.

\end{abstract}
\hrulefill

\section{Presentation}

ELCVIA is an international electronic journal on research and
applications in Computer  Vision and Image Analysis. The journal
has the following aims:

\begin{itemize}
\item To be a forum for discussion and exchange of theoretic
and methodological parts of Computer Vision and Image Analysis
research and their applications.

\item To be a dynamic and fast mean of communication between
researchers both in the revision of articles and in their publication.

\item To take advantage of Internet multimedia capacity.

\item To be a quality publication and therefore all articles will
be thoroughly revised. Moreover, the Journal aims to include
internationally recognized experts in
Computer Vision and Image Analysis in their Boards.
The objective is to include  ELCVIA in the {\em Journal Citation Index}.

\item To provide information on activities
(International Conferences, Workshops, PhDs, books etc) in
Computer Vision and Image Analysis.
\end{itemize}

The ELCVIA project started at the ICPR?2000 conference in Barcelona,
when a group  of researchers of international prestige expressed
their general interest in creating an electronic journal based on
new technologies, as its means of communication. This e-journal
would allow investigators to accelerate the process of keeping up
with the new advances in the field of Computer Vision and Image
Analysis without loosing quality and the rigour that scientific
articles require.

\section{Why an electronic journal?}

Internet has evolved into a means of communication and, from this
concept, many  electronic journals have emerged recently. These
e-journals have several advantages over paper-based publications:

\begin{itemize}
\item More compact storage which does not deteriorate.
\item Easy duplication without losses.
\item Possibility of printing just the articles you consider essential.
\end{itemize}

There are many paper-based journals that have also adopted Internet
as their form of publication. Nonetheless, we believe that they
restrict their services, since they are dependent on the paper format.
They only offer the possibility of obtaining an electronic copy of the
same article in paper format and don?t take advantage of the multimedia
capabilities that Internet gives.

ELCVIA is an electronic journal, without a paper based edition, that
tries to give authors the use of all the possibilities that Internet
offers for a better dissemination of their scientific research. The
main services ELCVIA intends to provide are:

\begin{itemize}
\item Attachment of multimedia files with the articles
\item Executable programs, scripts etc...
\item Animations (AVI and MPEG files).
\item Images(GIF, TIFF, JPEG) in colour.
\item Links to other references specified by the author.
\item Flexible size of articles and surveys.
\item Revision process over e-mail to make it faster and more
      agile without decreasing quality.
\item Search engine to make article tracking easier through keywords.
\item Personified mail to receive the articles of interest
specified by keywords, authors, affiliation, etc.
\item Low subscription prices. Institutional subscription by domain.
\item Draft publication.
\end{itemize}

\section{Scope}

The journal encourages research groups from all over the world
to send in their  contributions on computational aspects of Computer
Vision and Image Analysis and applications.

The authors can send three types of contributions:

\begin{description}
\item[SURVEY:] A contribution by invitation of the Editors
in Chief or Advisory Board editors, of 12-20 pages (recommended),
not revised by the Editorial Board and reviewers, to be accepted
for a regular issue of the journal.

\item[ARTICLE:] A regular contribution of 8-12 pages (recommended)
revised by the Editorial Board and reviewers, to be accepted for
a regular issue of the journal.

\item[DRAFT:] An article of 8-12 pages (recommended) free-submitted
in the web site with a forum which will include the remarks and
suggestions of everyone, subscriber or not, who will read this
draft. With these remarks, the author can decide to send to ELCVIA
a modified version of the draft as an article.

\end{description}

The different fields covered (but not limited to these) by
the journal are:

\begin{itemize}
\item Computer Vision and Image Analysis algorithms
\item Artificial Intelligence approaches in Computer Vision and
      Image Analysis
\item Pattern Recognition
\item Image segmentation
\item 3D reconstruction
\item Active vision
\item Tracking
\item Video and image sequence analysis
\item Vision applied to robot navigation
\item Computer Vision Architectures

\item Vision Applications:

\begin{itemize}
\item Document Analysis and OCR
\item Medical Image Analysis
\item Machine Vision
\item Robotics
\end{itemize}
\end{itemize}

\section{Periodicity}

\begin{itemize}
\item 6 issues per year.
\item The publication of the articles of a issue will be announced
      to subscribers via e-mail. Every time a new article is published,
      an E-mail will be sent announcing its publication. In order to
      avoid sending these e-mails very often, they will be sent every
      3 or 4 weeks.

\item Subscribers can select/decide on a set of keywords and the
      system can automatically send an e-mail with the reference.
\item The publication announcement sent by e-mail will include
      the authors and title of each article published and a link to
      each one in the journal Web site.
\item Access to article abstracts will be free and the complete papers will be
      free temporaly.
\item There will be two subscription types: individual and institutional,
      requiring user passwords and/or IP addresses. This will
      determine the access to articles.
\end{itemize}

\section{Contents}

Each number will include a survey and a variable number of articles
(5-10) that will  appear as soon as they are accepted by the referees:

\begin{description}
\item[SURVEYS:] We recommend 12-20 pages including figures. The surveys are
under invitation of the Editors in Chief or Advisory Board editors.

\item[ARTICLES:] We recommend 8-12 pages including figures. The Editorial
Board will determine the final length during the revision process.

The authors will send an electronic copy of the article in \LaTeX, Word,
PS or PDF format with multimedia files, and if it is accepted,
the publishers will make the conversion to other formats.

Each article shall have an HTML document with the abstract and
the article in  PostScript and PDF, as well as the multimedia files
(AVI, MPEG, GIF, etc) that complement the article.

\item[NEWS:] Activities related to Computer Vision and Image Analysis.
\begin{itemize}
\item Recent PhD abstracts.
\item Publication of Computer Vision related books.
\item Announcements (Call for articles, Call for participation).
\item Reports from conferences and meetings.
\item New interesting links.
\item Any other information that may be of interest to the community of
researchers in Computer Vision and Image Analysis.
\end{itemize}
\end{description}

\section{Boards}

ELCVIA seeks to join in their Boards recognized international specialists in Computer 
Vision and Image Analysis.

The ELCVIA Board will be a 5 level board:

\begin{description}
\item[Editors in Chief:] Responsible for strategic decisions.

\item[Advisory Board:] Together with the Editors in Chief,
they will be responsible for the journal quality and exempt from
executive tasks. They can propose special editions and invite other
researchers to contribute surveys. They propose and elect the Editorial
Board once every two years.

\item[Editorial Assistant:] Main co-ordinators of the journal. They
receive an article sent by the authors and forward it to one
Editorial Board member for review.

\item[Editorial Board:] They receive from the Ed. Assistants an
article to be reviewed and forward it to two reviewers.

\item[Reviewers:] Proposed for each separate article by an Editorial
Board member.
\end{description}

\section{Revision and discussion}

Two types of contributions are under consideration:

\begin{itemize}
\item {\bf Draft discussion:}

\begin{itemize}
\item The author sends a draft article (named draft), which is open
to everyone, opening a free-for-all forum and allowing the author
to receive comments and suggestions from the scientific community.
These suggestions are open to everyone, not only to the authors.
The draft will be open for revision for 2-3 months.

\item Draft discussion will not be blind, i.e. the authors of the 
articles  and their affiliation will be written in the draft. 

\item The subscribers will receive the announcement of new drafts
by e-mail in the same format as for normal articles (authors,
title and a link to the ELCVIA web site).
\end{itemize}

\item {\bf Article revision:}

\begin{itemize}
\item The author submits the article to the journal. The revision
process is as follows:

\begin{enumerate}

\item The Advisory Board editors together with the Editorial
Assistants choose an Editorial Board member named {\bf Associate Editor (AE)}
who will be responsible for the article revision.

\item The AE, then, chooses two reviewers (one of these may be the same AE) who
will review the article and then send their comments back to him.
He will then decide whether the article is appropriate for publication
and what length it should have.

\item The accepted article will be published in ELCVIA upon the
recommendation of the AE. For each published
paper will appear "Recommended for acceptance by $<$AE name$>$.

\end{enumerate}

\item Article revision will not be blind, i.e. the authors of the articles 
and their affiliation will be known to the AE, who handles the paper and 
the review.

\end{itemize}
\end{itemize}

\section{Special Issues}

\begin{itemize}

\item Special Issues may be proposed:

\begin{itemize}

\item By topic
\item By conference or Workshop (selection of the best articles of the congress)
\item By research groups
\item By country
\end{itemize}

\item In all cases, the articles of a special issue will be reviewed.

\end{itemize}

\section{Information to the Authors}

\begin{itemize}
\item Submission of an article or draft and the sending of the electronic files
is made by filling a Web form named {\em Submit an article or draft} located 
in the Author's Section of the ELCVIA Web site. The authors must give the 
following information:

\begin{enumerate}

\item Article information: Title, Authors, E-mail, Predefined Keywords 
or not (we suggest to use some Predefined Keywords to make easy the 
revision process), Article Type and Abstract.

\item The article file must be a PDF (recommended),  Postcript or Word file using our 
camera-ready format, offered in \LaTeX 2e, \LaTeX 2.09 or  Word'2000 formats. 

\item We encourage to the authors to send additional files (video, image of .exe) 
in standard formats (GIF, JPEG, AVI, MPEG) or in compressed (.ZIP) format to 
improve the discussion of the article or draft.

\end{enumerate}

\item When an article is accepted, ELCVIA Board will send to the authors the following 
information to make the final camera-ready version:

\begin{itemize}
\item Volume, number, initial and final pages, year.

\item Name of the Associate Editor of the article.

\item Received, revision (if it is necessary) and accepted dates.
\end{itemize}

This final camera-ready version must be send in PDF (recommended), Postcript 
or Word format.

\end{itemize}

\subsection{Style}

\noindent Camera-ready format.

\subsubsection{Sizes}

The size and margins of the document are shown in Table \ref{table:1}.

\begin{table}[h]
\centerline{
\begin{tabular}{|l||l|}
\hline
Page's size           & A4 (210 mm $\times$ 297 mm) \\ \hline \hline
Margins               &                             \\ \hline
Top                   & 2 cm                        \\
Left                  & 2 cm                        \\ \hline \hline
Text and figures area & 17 cm (width) $\times$ 25 cm (height) \\
\hline
\end{tabular}
}
\caption{Paper size and margins}
\label{table:1}
\end{table}

\subsubsection{Fonts}

The font type and size are shown in Table \ref{table:2}.

If the font ``Times" is not available another similar font
can be used, for example \LaTeX users can use the default
font {\em computer-modern-font}.

\begin{table}
\centerline{
\begin{tabular}{|l|l||l|}
\hline         
Title                      &  Times  16pt (bold)  \\
Section title              &  Times  14pt (bold)  \\
Subsection title           &  Times  12pt (bold)  \\
Subsubsection title        &  Times  11pt (bold)  \\
Abstract                   &  Times  10pt         \\
Body Text                  &  Times  11pt         \\
Figure and table captions  &  Times  11pt         \\
\hline
\end{tabular}
}
\caption{Fonts}
\label{table:2}
\end{table}

 
\subsubsection{Important}

\begin{itemize}

\item
      \LaTeX 2.09 users can use the style file
      {\tt elcvia.sty} and the example file {\tt elcvia209.tex}.

\item
      \LaTeX 2e users can use the style files
      {\tt elcvia.sty} and {\tt elcvia.cls} with the example file
      {\tt elcvia2e.tex}.

\item
     Word users can use the example file 
     {\tt elcvia.doc} and the macros file
     {\tt elcvia.dot}.

\item For any question related to the format please contact: {\sf elcvia@cvc.uab.es}

\end{itemize}

\noindent
{\Large \bf Note:}
This document has been formatted using the \LaTeX 2e style \cite{Les:94}
and its heading is:

\begin{verbatim}
\documentclass[twoside,a4paper,11pt]{article}
\usepackage{times}
\usepackage{elcvia}
\usepackage[english]{babel}
\end{verbatim}

If the font ``Times" is not available use this heading:

\begin{verbatim}
\documentclass[twoside,a4paper,11pt]{article}
\usepackage{elcvia}
\usepackage[english]{babel}
\end{verbatim}

{\tt article} corresponds to the standard \LaTeX files {\tt article.sty} and
{\tt article.cls}.

\section{Publisher:}

\begin{verbatim}
Computer Vision Center
Edifici O - Campus UAB
08193 Bellaterra (Barcelona) - Spain
e-mail: elcvia@cvc.uab.es
Tel: 34 - 93 - 581 18 28
Fax: 34 - 93 - 581 16 70
\end{verbatim}

\begin{thebibliography}{1}

%Book Reference
\bibitem{Les:94}
Leslie Lamport, {\em A Document Preparation System}, Addison-Wesley, 1994. 
{\em (Example of Book Reference)}

%Journal Reference
\bibitem{Phil:99}
I.T. Phillips, A.K. Chhabra, "Empirical performance evaluation of graphic
recognition systems", {\em IEEE Transactions on Pattern Analysis and Machine
Intelligence} 21(9):849-870, 1999.
{\em (Example of Journal Reference)}

%Proceedings Reference
\bibitem{KIK:86}
H. Kida, O. Iwaki, K. Kawada, "Document recognition system for office
automation", {\em IEEE Proc. of 8th Intl. Conf. on Pattern Recognition},
Paris, 1:446-448, 1986. 
{\em (Example of Proceedings Reference)}

\end{thebibliography}


\end{document}
